\color{unidarkgreen}{\subsection{Защита от непереключения фаз}}
\color{black}

\begin{enumerate}[label=\arabic{section}.\arabic{subsection}.\arabic*,  align=left, labelsep=4pt, leftmargin=0pt, itemindent=57pt]

\item
Защита от непереключения фаз (ЗНФ) предназначена для обнаружения расхождения полюсов
выключателей В1 ВН и В2 ВН с пофазным приводом при включении или отключении.   

\begin{figure}[H]
\centering
\includegraphics[width=0.88\textwidth,height=0.88\textheight,keepaspectratio]{\fbpath/761.3441 ЗНФ/img/ZNF.pdf}
\captionof{figure}{Функциональная схема ЗНФ В1 ВН и В2 ВН}\label{pict:ZNF}
\end{figure}

\item
Ввод в работу ЗНФ выключателя В1 ВН и В2 ВН осуществляется с помощью программных
переключателей \textbf{«ЗНФ В1 ВН»} и \textbf{«ЗНФ В2 ВН»}. Факт расхождения полюсов может устанавливаться одним из
следующих способов: 
\begin{itemize}
\item на основании внешнего сигнала от сборки блок-контактов выключателя;
\item функцией фиксации положения выключателя (ФПВ), которая пофазно контролирует состояние блок-контактов выключателей и в случае расхождения полюсов формирует сигналы \textit{«Внутр. пуск ЗНФ В1 ВН»} и \textit{«Внутр. пуск ЗНФ В2 ВН»}.
\end{itemize}

\item
В зависимости от используемого способа соответствующий сигнал должен назначаться на
конфигурируемый вход «Пуск ЗНФ В1 ВН» и «Пуск ЗНФ В2 ВН».  

\item
Действие на отключение выключателя В1 ВН и В2 ВН с запретом АПВ осуществляется с выдержками
времени \textit{«Тср ЗНФ В1 ВН»} и \textit{«Тср ЗНФ В2 ВН»}. Задаваемые выдержки времени предназначена для отстройки
от разновременности переключения блок-контактов выключателя.

\item
Предусмотрена возможность блокировки ЗНФ В1 ВН и В2 ВН от внешних сигналов \textit{«Блок. ЗНФ В1 ВН»} и
\textit{«Блок. ЗНФ В2 ВН»} на время, которое задаётся уставками \textbf{«Тдеблок ЗНФ В1 ВН»} и \textbf{«Тдеблок ЗНФ В2 ВН»}.

\item
Оперативный ввод/вывод ЗНФ В1 ВН и В2 ВН производится с помощью БУ (блоков управления
функциями) \textit{«Вывод ЗНФ В1 ВН»} и \textit{«Вывод ЗНФ В1 ВН»}.


\small
\setlength{\extrarowheight}{0.05cm} % высота ячеек в таблице
\begin{longtable}
{|>{\centering\arraybackslash}m{4.5cm}|>{\centering\arraybackslash}m{4.5cm}|>{\centering\arraybackslash}m{3cm}|>{\centering\arraybackslash}m{3.5cm}|}
\caption{Перечень уставочных параметров контроля положения выключателей\hfill\vspace{-0.5\baselineskip}}\label{znf:tbl1} \\ 
\hhline{|-|-|-|-|}
\rowcolor{gray!30}
Наименование уставки & Диапазон & Шаг &  Значение по умолчанию \\ 
\hhline{|-|-|-|-|}
\endfirsthead

\caption*{\hspace{3pt}\emph{Продолжение таблицы \ref{znf:tbl1}\hfill\vspace{-0.5\baselineskip}}} \\ % сделано по ГОСТ 2.105 п.6.8.7
\hhline{|-|-|-|-|}
\rowcolor{gray!30}
Наименование уставки & Диапазон & Шаг &  Значение по умолчанию \\  
\hhline{|-|-|-|-|}
\endhead
\hhline{|-|-|-|-|}

\endfoot
\endlastfoot

%%%%%%%%%%%%содержимое%%%%%%%%%%%%%%

\end{longtable}
\setlength{\extrarowheight}{0.0cm} % возврат параметра высота ячеек в таблице
\normalsize










































\end{enumerate} 